%%
%% Test: Standard LaTeX Package Conventions — v1.1
%% Stage: package conventions (stages 1-4 of the refactor)
%% 
%% Purpose: Verify that MatinBook now follows official LaTeX
%%          package conventions:
%% 
%%   Stage 1: \ProvidesPackage{name} in all .sty files
%%   Stage 2: \RequirePackage{name} in matinbook.cls
%%   Stage 3: Simple package names (no paths)
%%   Stage 4: Successful compilation
%% 
%% Run with:
%%   TEXINPUTS=<project>:<project>/tex/... xelatex stage-packages-01.tex
%% 

\documentclass{matinbook}

\title{تست استانداردهای پکیج — نسخه ۱.۱}
\author{تیم MatinBook}
\date{۱۴۰۵}

\begin{document}

\maketitle

\chapter{بررسی استانداردهای پکیج لاتک}
\label{chap:test-packages}

این سند، استانداردهای رسمی لاتک را در پروژه‌ی MatinBook
بررسی می‌کند.

%----------------------------------------
\section{مرحله ۱ — اعلان پکیج}
\label{sec:stage1}

\subsection{استفاده از \texttt{\textbackslash ProvidesPackage}}

در نسخه‌ی قبل، فایل‌های \verb|.sty| از \verb|\ProvidesFile| استفاده
می‌کردند که خلاف استاندارد لاتک است. در نسخه‌ی ۱.۱، همه‌ی ۲۴ فایل
به \verb|\ProvidesPackage| مهاجرت کردند.

\textbf{بررسی:} اگر این سند کامپایل شود، یعنی همه‌ی پکیج‌ها
به‌درستی اعلان شده‌اند.

\subsection{نام ساده بدون مسیر}

نام پکیج‌ها حالا ساده است (\verb|mb-engine|, \verb|mb-options|, ...)
بدون اسلش یا مسیر.

%----------------------------------------
\section{مرحله ۲ — بارگذاری پکیج}
\label{sec:stage2}

\subsection{استفاده از \texttt{\textbackslash RequirePackage}}

در \verb|matinbook.cls|، همه‌ی ماژول‌ها با \verb|\RequirePackage|
بارگذاری می‌شوند، نه \verb|\input|.

\textbf{بررسی:} هیچ خطای «requested document class» وجود ندارد.

%----------------------------------------
\section{مرحله ۳ — نام‌های ساده}
\label{sec:stage3}

\subsection{بدون هشدار «requested package»}

در نسخه‌ی قبل، هشدارهایی مثل این وجود داشت:

\begin{quote}
\verb|LaTeX Warning: You have requested package `tex/core/mb-engine',| \\
\verb|               but the package provides `mb-engine'.|
\end{quote}

\textbf{بررسی:} اگر در لاگ کامپایل چنین هشداری نیست، یعنی مرحله ۳
موفق بوده است.

\subsection{پکیج‌های تم}

پکیج‌های تم حالا با پیشوند \verb|mb-theme-| نام‌گذاری شده‌اند:
\begin{itemize}
    \item \verb|mb-theme-default|
    \item \verb|mb-theme-colors|
    \item \verb|mb-theme-fonts|
    \item \verb|mb-theme-cover|
\end{itemize}

%----------------------------------------
\section{مرحله ۴ — بررسی عملکردی}
\label{sec:stage4}

\subsection{محیط‌های علمی}

برای اطمینان از اینکه همه‌ی ماژول‌ها بارگذاری شده‌اند:

\begin{theorem}[قضیه نمونه]
    برای هر عدد حقیقی $x$، داریم $x^2 \geq 0$.
\end{theorem}

\begin{proof}
    چون مربع هر عدد حقیقی نامنفی است، نتیجه حاصل می‌شود.
\end{proof}

\begin{definition}[عدد اول]
    عدد طبیعی بزرگ‌تر از ۱ که تنها دو مقسوم‌علیه مثبت دارد،
    \emph{عدد اول} نامیده می‌شود.
\end{definition}

\subsection{ریاضیات}

\begin{equation}
    \label{eq:package-test}
    f(x) = \sum_{i=1}^{n} x_i^2
\end{equation}

مجموعه‌های اعداد: $\N \subset \Z \subset \Q \subset \R \subset \C$.

\subsection{کد}

\begin{latin}
\begin{minted}{python}
def greet(name):
    """Return a greeting."""
    return f"Hello, {name}!"
\end{minted}
\end{latin}

%----------------------------------------
\section{نتیجه‌گیری}
\label{sec:conclusion}

اگر این سند بدون خطا کامپایل شده باشد:

\begin{itemize}
    \item ✅ همه‌ی ۲۴ پکیج با \verb|\ProvidesPackage| اعلان شده‌اند
    \item ✅ \verb|matinbook.cls| از \verb|\RequirePackage| استفاده می‌کند
    \item ✅ نام پکیج‌ها ساده است (بدون مسیر)
    \item ✅ هشدار «requested package» وجود ندارد
    \item ✅ پکیج‌های تم با \verb|mb-theme-| نام‌گذاری شده‌اند
    \item ✅ محیط‌های علمی، ریاضی، و کد کار می‌کنند
\end{itemize}

\textbf{وضعیت تست:} قبول

\end{document}